% Руководство для агента. Для нового отчёта начните с ../../report.tex.
% Сборка примера из этой директории: TEXINPUTS=..: latexmk -lualatex bmstu-examples.tex
% В больших отчётах главы можно хранить отдельно и подключать через
% \chapter{Технологический раздел}
, как в ../../report.tex.
\documentclass{bmstu}
\usetikzlibrary{arrows.meta,positioning,shapes.geometric}
\addbibresource{biblio.bib}

\begin{document}

% Аргументы, помеченные как необязательные, могут быть пустыми. В таком случае соответствующее этому аргументу поле (например, ФИО консультанта) добавлено не будет.

% Титульная страница

% Отчет по лабораторной работе

% \makereporttitle
%     {Информатика и системы управления} % Название факультета
%     {Программное обеспечение ЭВМ и информационные технологии} % Название кафедры
%     {лабораторной работе №~1} % Название работы (в дат. падеже)
%     {Операционные системы} % Название курса (необязательный аргумент)
%     {Обработчик прерывания от системного таймера} % Тема работы
%     {} % Номер варианта (необязательный аргумент)
%     {Иванов~И.~И./ИУ7-51Б} % ФИО студента/номер группы (если авторов несколько, их необходимо разделить запятой)
%     {Петров~П.~П.} % ФИО преподавателя

% Расчетно-пояснительная записка к курсовой работе

\makecourseworktitle
    {Информатика и системы управления} % Название факультета
    {Программное обеспечение ЭВМ и информационные технологии} % Название кафедры
    {Практическое руководство по оформлению отчётов в LaTeX} % Тема работы
    {Иванов~И.~И./ИУ7-71Б} % ФИО студента/номер группы (если авторов несколько, их необходимо разделить запятой)
    {Петров~П.~П.} % ФИО научного руководителя
    {Сидоров~С.~С.,Кузнецов~К.~К.} % ФИО консультанта (необязательный аргумент; если консультантов несколько, их необходимо разделить запятой)

% Расчетно-пояснительная записка к научно-исследовательской работе

% \makeresearchtitle
%     {Информатика и системы управления} % Название факультета
%     {Программное обеспечение ЭВМ и информационные технологии} % Название кафедры
%     {Исследование методов генерации исходного кода} % Тема работы
%     {Иванов~И.~И./ИУ7-81Б,Сергеев~С.~С./ИУ7-82Б} % ФИО студента/номер группы (если авторов несколько, их необходимо разделить запятой)
%     {Петров~П.~П.} % ФИО научного руководителя
%     {} % ФИО консультанта (необязательный аргумент; если консультантов несколько, их необходимо разделить запятой)

% Расчетно-пояснительная записка к выпускной квалификационной работе

% \makethesistitle
%     {Информатика и системы управления} % Название факультета
%     {Программное обеспечение ЭВМ и информационные технологии} % Название кафедры
%     {Метод обнаружения спутников как точечных объектов на изображении} % Тема работы
%     {Иванов~И.~И./ИУ7-81Б} % ФИО студента/номер группы (если авторов несколько, их необходимо разделить запятой)
%     {Петров~П.~П.} % ФИО научного руководителя
%     {Сидоров~С.~С.} % ФИО консультанта (необязательный аргумент; если консультантов несколько, их необходимо разделить запятой)
%     {Кузнецов~К.~К.} % ФИО нормоконтролера

\begin{essay}{LaTeX, шаблон отчёта, библиография, таблицы, листинги, иллюстрации}
Объект работы~--- подготовка учебных отчётов в~классе \texttt{bmstu}.

Цель работы~--- дать агенту рабочие образцы оформления и~сборки отчёта.

Метод работы~--- проверка примеров в~собранном документе и~сверка правил оформления с~ГОСТ~7.32-2017.

Результат работы~--- руководство с~примерами вводных разделов, формул, листингов, таблиц, графиков и~приложений.
Руководство предназначено для подготовки лабораторных, курсовых и~исследовательских работ.
\end{essay}

\maketableofcontents

\begin{definitions}
    \definition{Листинг}{оформленный фрагмент исходного кода с~номером и~подписью}
    \definition{Перекрёстная ссылка}{ссылка внутри документа на~нумерованный элемент}
    \definition{Шапка таблицы}{строки с~заголовками столбцов}
\end{definitions}

\begin{abbreviations}
    \definition{ГОСТ}{государственный стандарт}
    \definition{НИР}{научно-исследовательская работа}
    \definition{CSV}{текстовый формат табличных данных с~разделителями}
    \definition{PDF}{формат переносимых документов}
    \definition{$n$}{число элементов множества}
\end{abbreviations}

\chapter*{ВВЕДЕНИЕ}
\addcontentsline{toc}{chapter}{ВВЕДЕНИЕ}

Руководство показывает, как агенту подготовить отчёт, подобрать команды класса \texttt{bmstu} и~проверить итоговый PDF.
Правила с~номерами пунктов относятся к~ГОСТ~7.32-2017 для отчётов о~НИР~\cite{gost732}.
Для учебной работы уточняйте требования кафедры и~методических указаний к~конкретному заданию.

\chapter{Подготовка отчёта}

\section{Структура и сборка}

Для нового отчёта возьмите \texttt{report.tex} из~корня шаблона.
Выберите одну титульную команду по~виду работы: \verb|\makereporttitle|, \verb|\makecourseworktitle|, \verb|\makeresearchtitle| или~\verb|\makethesistitle|.
В~исходнике этого руководства показаны их аргументы.
Данные для титульного листа берите из~задания.
Отсутствующие обязательные реквизиты уточняйте у~автора, а~необязательные аргументы можно оставить пустыми.

\textbf{Курсовые работы, курсовые проекты, ВКР и~аналогичные расчётно-пояснительные записки.}
Следуйте структуре \texttt{report.tex}.
Сохраните вводные разделы, в~том числе реферат, содержание, термины и~определения, перечень сокращений и~обозначений и~введение.
Основную часть оформляйте в~порядке шаблона: аналитический, конструкторский, технологический и~исследовательский разделы.
Далее идут заключение, список источников и~приложения при наличии.

\textbf{Отчёты о~НИР.}
Включайте реферат, содержание, термины и~определения, сокращения и~обозначения, введение, заключение, список источников и~приложения при наличии.
Названия и~состав глав основной части определяйте по~задачам исследования.
Аналитический, конструкторский, технологический и~исследовательский разделы из~шаблона для НИР сами по~себе не~обязательны.

\textbf{Лабораторные работы.}
По~умолчанию обязателен только титульный лист.
Остальную структуру подбирайте по~заданию и~указаниям преподавателя.
Например, вместо «Введения» можно сделать «Задание», а~далее оформить ход работы и~результаты.
Добавляйте содержание, реферат и~перечни, когда они нужны по~заданию.

Имена \texttt{.tex}-файлов можно менять по~содержанию.
После переименования обновите вызовы \verb|\include| и~\verb|\input|.
Если меняется заголовок ненумерованного раздела, обновите и~запись в~содержании, если она есть.
Например, в~лабораторной можно переименовать \texttt{10-introduction.tex} в~\texttt{10-task.tex} и~заменить заголовок на~«ЗАДАНИЕ».
Для работы с~заданной структурой переименование файла сохраняет требуемое название раздела в~PDF.

Основные главы удобно хранить в~отдельных \texttt{.tex}-файлах и~подключать через \verb|\include|.
Небольшие фрагменты без принудительного начала страницы подключайте через \verb|\input|.
Рисунки для \verb|\includeimage| храните в~\texttt{inc/img/}, исходники для \verb|\includelisting|~--- в~\texttt{inc/lst/}, библиографию~--- в~\texttt{biblio.bib}.

В~обычном тексте \texttt{.tex}-файлов соблюдайте правило «одно предложение~--- одна строка исходника».
Начинайте каждое новое предложение с~новой строки и~оставляйте длинные предложения целиком на~одной строке.
Пустую строку оставляйте между абзацами.
LaTeX сам переносит текст в~PDF, поэтому ручные переносы внутри предложения и~команда \verb|\\| для этого не~нужны.

Собирайте из~папки главного \texttt{.tex}-файла.
Для этого руководства используйте команду ниже.
Для \texttt{report.tex} из~корня шаблона задайте \texttt{TEXINPUTS=bmstu:} и~передайте \texttt{report.tex} вместо имени примера.
\texttt{latexmk} запускает повторные проходы и~Biber по~мере необходимости.

\begin{listing}{bash}{build-guide}{Сборка руководства из каталога examples}
TEXINPUTS=..: latexmk -lualatex bmstu-examples.tex
\end{listing}

После правок пересоберите PDF.
Проверьте предупреждения о~потерянных и~повторных метках, номера страниц в~содержании, библиографию и~места разрыва таблиц.
Сверьте также, что подписи и~формулы целиком помещаются между полями.

\section{Вводные разделы}

Для реферата используйте окружение \texttt{essay}.
По~ГОСТ~7.32-2017, пп.~5.3 и~6.12, реферат отчёта о~НИР включает сведения об~объёме, перечень из~5--15 ключевых слов и~текст.
Ключевые слова дают в~именительном падеже прописными буквами, через запятую и~без точки в~конце.
В~тексте последовательно указывают объект, цель, методы, результаты и~применение.
Сведения, которых в~работе нет, пропускают.
После сборки сверяйте автоматические счётчики реферата с~фактическим содержимым.

Содержание выводит \verb|\maketableofcontents|.
Ненумерованные введение и~заключение добавляйте в~содержание явно через \verb|\addcontentsline|.
Для реферата, определений, сокращений и~библиографии это уже делают команды класса.
ГОСТ требует, чтобы названия в~содержании соответствовали заголовкам, а~номера страниц были справа и~соединялись с~названиями отточием (пп.~5.4 и~6.13).

В~окружениях \texttt{definitions} и~\texttt{abbreviations} используйте \verb|\definition{термин}{пояснение}|.
Перечисляйте только обозначения и~термины из~самой работы.
Упорядочивайте записи по~алфавиту, разделяйте название и~пояснение тире, оставляйте строки без конечного знака препинания (пп.~6.14 и~6.15).

\section{Рисунки}

Для готового изображения используйте \verb|\includeimage|.
Укажите имя без расширения, режим \texttt{f}, позицию, ширину и~подпись.
Класс ищет файл в~\texttt{inc/img/} и~создаёт метку \texttt{img:tux} для примера на~рисунке~\ref{img:tux}.
Режим \texttt{w} обтекает изображение текстом.
Для него позиция \texttt{r} означает правую сторону, а~\texttt{l}~--- левую.

\includeimage{tux}{f}{htbp}{0.20\linewidth}{Символ Linux}

Ссылайтесь на каждый рисунок перед его появлением.
Помещайте рисунок после первого упоминания или на~следующей странице.
Подпись ставят под рисунком по~центру и~пишут с~прописной буквы без точки в~конце (пп.~6.5.1, 6.5.7 и~6.5.8).

Для составного рисунка используйте \verb|\includeimages|.
Каждая группа имён в~первом аргументе задаёт отдельный ряд.
Далее идут общая подпись и~имя метки без префикса \texttt{img:}.
На~рисунке~\ref{img:linux-mascots} сверху расположены Tux и~Tuz, а~снизу по~центру повторён Tux.
Буквы идут по~порядку через все ряды.
Если сетка целиком не~помещается на~страницу, разбейте её на~несколько рисунков.

\includeimages{{tux,tuz},{tux}}{Талисманы Linux в двух рядах}{linux-mascots}

\section{Формулы}

Для отдельной формулы используйте \texttt{equation}, для нескольких строк с~одним номером~--- \texttt{aligned} внутри \texttt{equation}.
Класс нумерует формулы сквозным счётом по всему документу.
Ссылку на~номер оформляйте через \verb|\eqref|.
Например, число сочетаний вычисляется по~формуле~\eqref{eq:combinations}.

\begin{equation}\label{eq:combinations}
    \binom{n}{k} = \frac{n!}{k!(n-k)!},
\end{equation}

где $n$~--- число элементов множества,\\
$k$~--- число выбираемых элементов.

По~п.~6.8 поясняйте новые обозначения сразу под формулой в~порядке их появления.
Каждое пояснение начинайте с~новой строки.
После «где» двоеточие не~ставят.
Для системы условий используйте \texttt{cases}, как в~формуле~\eqref{eq:threshold}.

\begin{equation}\label{eq:threshold}
    f(x)=
    \begin{cases}
        0, & x < 0,\\
        x, & x \geq 0.
    \end{cases}
\end{equation}

\begin{equation}\label{eq:aligned-example}
    \begin{aligned}
        S_2 &= 1 + 2,\\
        S_3 &= S_2 + 3.
    \end{aligned}
\end{equation}

В~формуле~\eqref{eq:aligned-example} символ \verb|&| задаёт выравнивание по~знаку равенства.
Для ненумерованной выкладки используйте \texttt{equation*}, для матрицы в~такой выкладке~--- \texttt{pmatrix}.

\chapter{Оформление результатов}

\section{Листинги}

Для файла используйте \verb|\includelisting| и~передавайте имя внутри \texttt{inc/lst/}.
Расширение можно опустить, если оно определяется однозначно.
Метка создаётся из~имени файла без расширения.
Если один файл нужно вставить несколькими диапазонами, передайте разным вызовам разные суффиксы в~последнем необязательном аргументе.
Для кода прямо в~\texttt{.tex} задайте язык, имя метки и~подпись в~окружении \texttt{listing}.
Отдельное имя функции в~абзаце оформляйте через \verb|\mintinline|.

Листинги~\ref{lst:greeting-header} и~\ref{lst:greeting-body} показывают два последовательных фрагмента файла \texttt{greeting.c}.
Параметры \texttt{firstline} и~\texttt{lastline} задают диапазон строк.
У~каждого фрагмента свой номер и~своя метка.

\includelisting[firstline=1,lastline=2]{c}{greeting}
    {Подключение заголовочного файла}[header]

\includelisting[firstline=3,lastline=7]{c}{greeting}
    {Функция \texttt{print\_greeting}}[body]

В~листинге~\ref{lst:call-greeting} показан вызов этой функции.

% В окружении код пишется прямо в .tex, а имя задаёт метку lst:call-greeting.
\begin{listing}{c}{call-greeting}{Вызов функции приветствия}
print_greeting("Ada");
\end{listing}

Для коротких и~длинных листингов используйте одну и~ту же команду или окружение.
Если код целиком не~помещается на~странице, класс переносит его автоматически и~над каждым следующим фрагментом выводит «Продолжение листинга» с~тем же номером.
Метка остаётся общей.
Пример приведён в~листинге~\ref{lst:report_stats} приложения~\ref{app:source}.
Требования к~оформлению кода сверяйте с~указаниями кафедры.

\section{Таблицы}

В~таблице~\ref{tab:autofit} приведены тестовые случаи с~короткими значениями и~описаниями разной длины.

Для таблиц используйте \texttt{longtblr}.
Класс задаёт полную ширину, сетку и~подписи слева.
В~параметрах таблицы укажите подпись, метку и~состав столбцов.
Параметр \texttt{rowhead=2} повторяет две строки шапки, а~\verb|\SetCell| с~параметром \texttt{r} или~\texttt{c} объединяет ячейки по~вертикали или~горизонтали.
Параметр \texttt{m} центрирует содержимое ячейки по~высоте.

По~п.~6.6 на~каждую таблицу нужна ссылка.
Название ставят слева над первой частью с~прописной буквы и~без точки.
Для продолжения указывают тот же номер таблицы.
Заголовки граф центрируют и~пишут без точки в~конце.
Заголовки строк выравнивают слева.
Разрыв страницы возможен между строками, поэтому проверяйте высокие ячейки и~группы объединённых строк в~PDF.

Для нескольких сносок объявите пояснения через \verb|note{1}={текст}|, \verb|note{2}={текст}| и~так далее.
Номера задаются вручную.
Вставляйте метку через \verb|\TblrNote{1}| в~каждой ячейке, к~которой относится первое пояснение.
Текст сноски под таблицей появится один раз.

% X[-1,c,m] учитывает ширину содержимого и центрирует его по обеим осям.
% X[l,m] заполняет оставшееся место, переносит текст и центрирует его по высоте.
\begin{longtblr}[
    caption={Проверка автоподбора ширины столбцов},
    label={tab:autofit},
    note{1}={Проверка минимального входного значения.},
    note{2}={Повторите опыт с~другим набором входных данных.},
]{
    colspec={X[-1,c,m] X[l,m] X[-1,c,m]},
    rowhead=2,
    row{1,2}={c},
}
    \SetCell[r=2]{c,m} № & \SetCell[c=2]{c,m} Проверка & \\
      & Описание & Результат \\
    1 & Пустая строка. & Пройдено\TblrNote{1} \\
    2 & В начале и в конце строки есть пробелы. В отчёте важно отличать пробелы, оставленные пользователем, от отступов, которые движок добавляет при вёрстке. & Проверить\TblrNote{2} \\
    3 & Одна буква «Я». & Пройдено\TblrNote{1} \\
    4 & «Значение есть?» --- «Да, но в скобках (вторая попытка)». Кавычки, вопросительный знак и скобки проверяются в одной ячейке. & Пройдено \\
    5 & Длинное слово «высокопроизводительный» стоит рядом с короткими словами, чтобы переносу пришлось выбирать между пробелом и допустимой точкой разрыва внутри слова. & Пройдено \\
    6 & Смешанная запись: русский текст, английское название \texttt{report cache} и число \(12{,}75\). & Пройдено \\
    7 & В строке подряд идут три коротких сообщения. Первое сообщает о запуске, второе показывает результат, третье поясняет, почему результат следует перепроверить на другом наборе данных. & Проверить\TblrNote{2} \\
    8 & Только число \(0\). & Пройдено \\
    9 & После обновления конфигурации приложение повторно обработало входные данные. Для проверки нужно сравнить две записи журнала и убедиться, что совпадают идентификатор, время и код завершения. & Проверить \\
    10 & Путь к файлу указан через \texttt{report / examples / output.pdf}. Здесь пробелы вокруг разделителей оставляют движку места для переноса. & Пройдено \\
    11 & Два предложения с очень разной длиной. Второе специально занимает больше одной строки, чтобы показать, как высота ячейки увеличивается вместе с объёмом текста. & Пройдено \\
    12 & Пропуск значения обозначен тире. & Пройдено \\
    13 & В тексте встречаются знак процента \(15\%\), знак равенства и единицы измерения. Короткое значение результата должно оставаться по центру своей ячейки. & Пройдено \\
    14 & Условная запись из журнала: задача запущена вручную в режиме проверки, обработка завершена успешно, итоговый файл сохранён рядом с исходными данными. & Пройдено \\
    15 & Необязательное поле оставлено пустым, поэтому для него применяется настройка по умолчанию. В журнале отдельно фиксируется, какое значение выбрала программа. & Проверить \\
    16 & Последняя строка подтверждает, что после разрыва страницы сохраняются те же ширины столбцов и та же шапка. & Пройдено \\
\end{longtblr}

\section{График по результатам измерений}

В~таблице~\ref{tab:timings} и~на рисунке~\ref{img:measurements} приведены условные результаты трёх опытов\footnote{В~этом примере каждый опыт соответствует одному запуску программы.}.
Значения заданы прямо в~исходнике строками таблицы и~координатами \texttt{pgfplots}.

Для пояснения внизу страницы используйте \verb|\footnote{текст}|.
По~п.~6.7.4 знак ставят надстрочно сразу после поясняемого фрагмента, без пробела.
Для пояснения под самой таблицей задайте \verb|note{*}={текст}| в~параметрах \texttt{longtblr}, а~в~ячейке поставьте \verb|\TblrNote{*}|.
Такая заметка относится только к~таблице.

Класс задаёт чёрно-белый стиль графиков.
При сравнении нескольких рядов различайте их формой маркеров и~типом линии.
Для графиков из~внешних файлов выбирайте монохромное оформление при экспорте и~проверяйте, что ряды различимы на~печати.

\begin{longtblr}[
    caption={Время выполнения по опытам},
    label={tab:timings},
    note{*}={Время округлено до целых миллисекунд.},
]{colspec={X[c,m] X[c,m]},rowhead=1}
    Номер опыта & Время, мс\TblrNote{*} \\
    1 & 12 \\
    2 & 18 \\
    3 & 25 \\
\end{longtblr}

\begin{figure}[htbp]
    \centering
    \begin{tikzpicture}
        \begin{axis}[
            width=0.75\linewidth,
            xlabel={Номер опыта},
            ylabel={Время, мс},
            xtick={1,2,3},
            grid=major,
        ]
            \addplot coordinates {(1,12) (2,18) (3,25)};
        \end{axis}
    \end{tikzpicture}
    \caption{Время выполнения по опытам}
    \label{img:measurements}
\end{figure}

Для столбчатых диаграмм используйте стиль \texttt{bmstu bars}.
Он чередует чёрно-белые узоры и~серую заливку.
На~рисунке~\ref{img:measurement-bars} первый ряд обозначен диагональными линиями, второй~--- точками.
Те же узоры показаны в~легенде.

\begin{figure}[htbp]
    \centering
    \begin{tikzpicture}
        \begin{axis}[
            bmstu bars,
            width=0.75\linewidth,
            height=6cm,
            bar width=9pt,
            xlabel={Номер опыта},
            ylabel={Время, мс},
            xtick={1,2,3},
            ymin=0,
            ymax=32,
            enlarge x limits=0.25,
            legend style={at={(0.5,1.03)},anchor=south,legend columns=2},
        ]
            \addplot coordinates {(1,12) (2,18) (3,25)};
            \addplot coordinates {(1,9) (2,14) (3,20)};
            \legend{Исходный режим, Изменённый режим}
        \end{axis}
    \end{tikzpicture}
    \caption{Сравнение времени выполнения в двух режимах}
    \label{img:measurement-bars}
\end{figure}

\section{Блок-схема на альбомной странице}

Окружайте \texttt{landscape} только широкий рисунок или таблицу с~подписью.
Заголовок раздела и~поясняющий текст оставляйте на~странице книжной ориентации.
На~рисунке~\ref{img:workflow} показан порядок проверки после сборки.
Схема нарисована в~TikZ.
Для диаграмм и~схем в~отчёте о~НИР учитывайте требования ЕСКД (п.~6.5.2).

\begin{landscape}
\begin{figure}[H]
    \centering
    \begin{tikzpicture}[
        node distance=24mm,
        block/.style={draw, rounded corners, align=center, minimum height=12mm, inner sep=6pt},
        link/.style={-{Stealth}},
    ]
        \node[block] (build) {Собрать PDF};
        \node[draw, diamond, aspect=2, right=of build] (check) {Сборка прошла?};
        \node[block, right=of check] (review) {Проверить PDF};
        \node[block, below=of check] (fix) {Исправить ошибку};
        \draw[link] (build) -- (check);
        \draw[link] (check) -- node[above]{да} (review);
        \draw[link] (check) -- node[right]{нет} (fix);
        \draw[link] (fix) -| (build);
    \end{tikzpicture}
    \caption{Порядок проверки сборки}
    \label{img:workflow}
\end{figure}
\end{landscape}

\section{Списки и источники}

В~перечислении, которое продолжает вводную фразу, начинайте пункты со~строчной буквы.
По~ГОСТ~7.32-2017, п.~6.4.6, простой список разделяют запятыми, сложный~--- точками с~запятой.
Последний пункт завершает предложение точкой.
Для перечня по~умолчанию ставят тире.
Если на~пункт нужно сослаться, используют строчные русские буквы со~скобкой, пропуская ё, з, й, о, ч, ъ, ы, ь.
Для известного числа пунктов допускаются цифры со~скобкой.
Класс задаёт маркеры, а~регистр и~знаки препинания нужно задать в~тексте.

Комплект отчёта включает:
\begin{itemize}[labelindent=\parindent,leftmargin=*]
    \item исходник,
    \item данные,
    \item собранный PDF.
\end{itemize}

После правок выполните три проверки:
\begin{enumerate}[labelindent=\parindent,leftmargin=*]
    \item пересоберите PDF и~проверьте журнал сборки;
    \item откройте изменённые страницы и~сверьте разбиение таблиц;
    \item проверьте номера элементов и~переходы по~ссылкам.
\end{enumerate}

Если в~тексте нужны ссылки на~отдельные пункты, можно выделить два условия готовности:
\begin{enumerate}[label=\asbuk*),labelindent=\parindent,leftmargin=*]
    \item все исходники и~используемые данные доступны;
    \item PDF успешно собирается и~доступен для проверки.
\end{enumerate}

Условие~б) проверяют после выполнения условия~а).

\begin{samepage}
Вложенные списки оформляйте только там, где у~пункта действительно есть подпункты.
Например, проверка отчёта включает:
\begin{enumerate}
    \item сборку документа;
    \item просмотр страниц с~результатами, включая:
    \begin{itemize}
        \item подписи рисунков,
        \item повтор шапки таблицы.
    \end{itemize}
\end{enumerate}
\end{samepage}

Ключ записи в~файле \texttt{biblio.bib} передаётся команде \verb|\cite|.
Например, описание пакета для оформления листингов приведено в~документации~\cite{minted}.

Список выводите через \verb|\makebibliography|.
По~пп.~6.9 и~6.16 ссылки оформляют номерами в~квадратных скобках, а~источники приводят в~порядке первого цитирования.
Заполняйте карточки только по~проверенным данным источника и~проверяйте после сборки, что все ключи разрешились.

\section{Приложения и проверка оформления}

Приложения создавайте внутри \texttt{appendices}, по~одной команде \verb|\chapter| на~приложение.
Ссылайтесь на~них в~порядке появления.
В~приложении~\ref{app:source} показан полный исходник, а~в~приложении~\ref{app:pdf}~--- вставка выбранной страницы готового PDF через \verb|\includepdf|.
Для вставленных страниц проверьте масштаб, ориентацию и~сквозную нумерацию (п.~6.17).

Для отчёта о~НИР сверяйте поля и~нумерацию с~пп.~6.1 и~6.3.
Учебный класс использует правое поле 10~мм, а~в~ГОСТ для отчёта о~НИР предусмотрено 15~мм.
Выбирайте параметры по~требованиям конкретной работы.
Номер титульного листа учитывают в~сквозной нумерации и~на~самом листе не~печатают.
В~заголовках разделов не~ставьте точку в~конце.
Класс задаёт абзацный отступ 1,25~см для первой строки нумерованного заголовка и отключает переносы слов в заголовках.
Вертикальные интервалы настроены по пропорциям схемы оформления РПЗ ВКР МГТУ.

\chapter*{ЗАКЛЮЧЕНИЕ}
\addcontentsline{toc}{chapter}{ЗАКЛЮЧЕНИЕ}

Перед выдачей отчёта проверьте итоговый PDF по~заданию и~требованиям кафедры.
Убедитесь, что в~документе нет учебных имён и~условных данных из~этого руководства, а~все ссылки ведут к~нужным элементам.

\makebibliography

\begin{appendices}
    \chapter{Исходный код программы}
    \label{app:source}
    \includelisting{c}{main}
        {Исходный файл \texttt{main.c}}

    Листинг~\ref{lst:report_stats} содержит программу для обработки CSV с~результатами измерений. Код продолжается на~следующих страницах.

    \includelisting{python}{report_stats}{Расчёт статистики по результатам измерений}

    \chapter{Вставка страницы готового PDF}
    \label{app:pdf}
    Для демонстрации вставлена первая страница файла с~эмблемой из~поставки класса. В~отчёте этим способом можно приложить лист задания или другой предоставленный документ.
    \includepdf[pages={1},scale=0.45,pagecommand={\thispagestyle{plain}}]{../bmstu-logo.pdf}
\end{appendices}

\end{document}
